% Sources: README "Model Physics" (intro); results pending
% Draft of the model part only; the results part waits for the full-range fits.

We present a model of the broadband spectrum, from the optical to the X-rays, of a supercritical accretion disc with an optically thick wind around a black hole, or around a neutron star with a magnetosphere.
Inside the spherization radius the disc loses mass to a wind launched along straight streamlines, and the accretion rate falls towards the centre.
Inside the trapping radius advection makes the effective temperature fall as $R^{-1/2}$.
The wind photosphere is found from the streamline density, and its temperature from the advection of photons by the wind together with the energy that the disc releases but does not show directly.
For a neutron star, the matter inside the magnetospheric radius radiates from a sphere or from the surface of the last closed dipole field line.
The observed spectrum is the sum of blackbodies over the visible surface at the inclination of the observer, with occultation by the wind and the disc.
% OUTLINE: Fits to the SEDs of the selected ULXs: method (MCMC) and main results.
